%% titre: Quiz 4 — Conditions d'optimalité
%% matiere: Optimisation
%% niveau: Licence 3
%% annee: 2026-2027
%% ordre: 4
%% description: Conditions nécessaires et suffisantes, points critiques, cas convexe, optimisation sur un convexe.

\begin{question}
Si $x^*$ est un minimum local d'une fonction $f$ différentiable sur $\R^n$, alors :
\begin{reponses}
  \bonne $\nabla f(x^*)=0$
  \mauvaise $\nabla^2f(x^*)$ est définie positive
  \mauvaise $f$ est convexe
  \mauvaise $\nabla f(x^*)\neq0$
\end{reponses}
\end{question}

\begin{question}
Si $f$ est de classe $\mathcal C^2$ et $x^*$ est un minimum local, alors $\nabla^2f(x^*)$ est :
\begin{reponses}
  \mauvaise définie positive
  \bonne semi-définie positive
  \mauvaise définie négative
  \mauvaise nulle
\end{reponses}
\begin{explication}
C'est la condition nécessaire d'ordre $2$. Elle est seulement \emph{semi}-définie positive : $x^4$ a un minimum en $0$ avec $f''(0)=0$.
\end{explication}
\end{question}

\begin{question}
Soit $f\in\mathcal C^2$, $\nabla f(x^*)=0$ et $\nabla^2f(x^*)$ définie positive. Alors $x^*$ est :
\begin{reponses}
  \mauvaise un maximum local
  \mauvaise un point selle
  \bonne un minimum local strict
  \mauvaise nécessairement un minimum global
\end{reponses}
\begin{explication}
C'est la condition suffisante d'ordre $2$ : elle donne un minimum \emph{local} strict. Pour un minimum global, il faut en plus la convexité.
\end{explication}
\end{question}

\begin{question}
Pour $f(x)=x^3$ en $x=0$ :
\begin{reponses}
  \mauvaise $0$ est un minimum local
  \mauvaise $0$ est un maximum local
  \mauvaise $0$ n'est pas un point critique
  \bonne $0$ est un point critique qui n'est ni un minimum ni un maximum local
\end{reponses}
\begin{explication}
$f'(0)=0$ et $f''(0)=0$, mais $f$ change de signe en $0$ : $\nabla f=0$ n'est pas suffisant.
\end{explication}
\end{question}

\begin{question}
Pour $f(x)=x^4$ en $x=0$ :
\begin{reponses}
  \mauvaise ce n'est pas un minimum, car $f''(0)=0$
  \bonne c'est un minimum global bien que $f''(0)=0$
  \mauvaise c'est un point selle
  \mauvaise c'est un maximum local
\end{reponses}
\begin{explication}
$x^4\ge0=f(0)$ pour tout $x$. La condition « hessienne définie positive » est suffisante, pas nécessaire.
\end{explication}
\end{question}

\begin{question}
En un point critique $x$, la hessienne a une valeur propre strictement positive et une valeur propre strictement négative. Alors $x$ est :
\begin{reponses}
  \mauvaise un minimum local strict
  \mauvaise un maximum local strict
  \bonne un point selle (ni minimum, ni maximum)
  \mauvaise un point où l'on ne peut rien dire
\end{reponses}
\end{question}

\begin{question}
Le point $(0,0)$ pour $f(x,y)=x^2-y^2$ est :
\begin{reponses}
  \mauvaise un minimum local
  \mauvaise un maximum local
  \mauvaise pas un point critique
  \bonne un point selle
\end{reponses}
\begin{explication}
$\nabla f=(2x,-2y)$ s'annule en $(0,0)$, et $\nabla^2f=\operatorname{diag}(2,-2)$ a deux valeurs propres de signes opposés.
\end{explication}
\end{question}

\begin{question}
En un point critique, la hessienne est $\begin{pmatrix}2&0\\0&0\end{pmatrix}$. Que conclut-on avec les conditions du second ordre ?
\begin{reponses}
  \mauvaise un minimum local strict
  \mauvaise un point selle
  \bonne le test est inconclusif : il faut étudier $f$ directement
  \mauvaise un maximum local
\end{reponses}
\begin{explication}
La hessienne est semi-définie positive avec une valeur propre nulle : par exemple $x^2+y^4$ a un minimum en $(0,0)$ alors que $x^2+y^3$ n'en a pas.
\end{explication}
\end{question}

\begin{question}
Si $f$ est convexe et différentiable sur $\R^n$, alors $x^*$ est un minimum global si et seulement si :
\begin{reponses}
  \bonne $\nabla f(x^*)=0$
  \mauvaise $\nabla^2f(x^*)\succ0$
  \mauvaise $f(x^*)=0$
  \mauvaise $\nabla f(x^*)\neq0$
\end{reponses}
\end{question}

\begin{question}
Soit $U$ un convexe et $x\in U$ un minimum local de $f$ différentiable sur $U$. Alors, pour tout $y\in U$ :
\begin{reponses}
  \mauvaise $\langle\nabla f(x),y-x\rangle<0$
  \mauvaise $\langle\nabla f(x),y-x\rangle\le0$
  \bonne $\langle\nabla f(x),y-x\rangle\ge0$
  \mauvaise $\langle\nabla f(x),y-x\rangle=0$
\end{reponses}
\begin{explication}
Aucune direction admissible ne fait décroître $f$ : $-\nabla f(x)$ ne pointe vers aucun point de $U$. L'égalité n'a lieu que lorsque $U$ est un sous-espace affine.
\end{explication}
\end{question}

\begin{question}
Soit $U=x+F$ un sous-espace affine et $x$ un minimum local de $f$ sur $U$. Alors :
\begin{reponses}
  \mauvaise $\nabla f(x)\in F$
  \bonne $\nabla f(x)\in F^{\perp}$
  \mauvaise $\nabla f(x)=0$ nécessairement
  \mauvaise $\nabla f(x)\notin F^{\perp}$
\end{reponses}
\begin{explication}
On a $\langle\nabla f(x),h\rangle=0$ pour tout $h\in F$. Le gradient n'est pas forcément nul : il est seulement orthogonal aux directions admissibles.
\end{explication}
\end{question}
